\scanpage{12}
- \mspage{9} -
\uline{et pour différentielle invariante} $\pi=dx/2y$.

De plus, sous ces hypothèses, la donnée d'une forme différentielle invariante suffit à rigidifier une courbe.

4.\quad \uline{Cas où} 2 \uline{est nilpotent et} $c_4$ \uline{inversible}.

Lorsque $2=0$, la forme différentielle ne suffit plus jamais à rigidifier une courbe de genre un. Lorsque $c_4$ est inversible et 2 nilpotent, donc $a_1$ inversible, en changeant le choix de $\pi,x,y$ on peut trouver une équation
\[
(4.1)\qquad y^2+xy=x^3+a_2x^2+a_6\qquad\pi=\frac{dx}{x+2y}.
\]
Les autres choix de $x$ et $y$ donnant lieu à une telle équation s'obtiennent en faisant, dans (1.2)
\[
(4.2)\qquad r=t=0\qquad u=1+2s
\]
et on aura
\[
\begin{aligned}
(4.3)\qquad u^2a'_2&=a_2-s-s^2\\
u^6a'_6&=a_6
\end{aligned}
\]
et \uline{modulo 2}
\[
(4.4)\qquad c_4=1\qquad c_6=1\qquad\Delta=a_6.
\]

5.\quad \uline{Cas où} $c_4=0$, $p=2$ \uline{ou} $p=3$

a) $p=3$

En caractéristique 3, lorsque l'équation est mise sour forme
\margin{11}
